
\section{EviStateBench Benchmark Design}
\label{sec:benchmark}

\subsection{Benchmark Package and Artifact Boundary}

\bench\ is a benchmark package with three components: public artifacts, hidden evaluator artifacts, and a shared evaluation protocol. The public artifacts contain task specifications, StateObservation streams, query sets, metadata, manifests, validation reports, and build commands. The hidden evaluator artifacts contain hidden state timelines and hidden answer sets. Systems under test read only the public package and submit QueryAnswer predictions. The evaluator reads the predictions and hidden answers and computes the reported metrics.
This boundary is central to the benchmark design. The public artifact exposes information that a deployed state-maintenance system could plausibly observe: task descriptions, observation records, query specifications, and metadata needed for parsing and normalization. It does not expose hidden timeline events, answer sets, or perturbation labels that reveal which observations were injected, removed, delayed, or contradicted. Thus, systems are evaluated on temporal task-state maintenance from public evidence, not on access to simulator truth.
The same boundary also makes the benchmark representation-agnostic. A system may implement a retrieval memory, temporal log, SQL scan, rule engine, materialized view, learned memory module, or hybrid architecture. It is not required to expose its internal representation. All systems are compared through the same QueryAnswer schema and the same hidden-oracle evaluator.

\engine\ is evaluated by \bench; it does not generate the ground truth used to judge other systems. The oracle is constructed from simulator truth and task specifications through the benchmark generation pipeline. This separation avoids making the reference engine part of the labeling process and allows \engine\ to be compared with simpler or future systems under the same public inputs and hidden answers.
The package supports three independent roles. A \emph{benchmark builder} regenerates public and hidden artifacts from pipeline manifests. A \emph{system developer} downloads the public package, implements an answerer, validates prediction format, and emits predictions. An \emph{evaluator} scores predictions against hidden answers. This separation is important for artifact review and public leaderboards: system developers can test parsers, run baselines, and validate output schemas without access to hidden timelines or answer files.
The public-artifact validation step checks that required manifests exist, query identifiers are well formed, observation streams are readable, task specifications align with query metadata, and answer leakage is absent from public files. In the reported artifact, validation covers 570 task specifications, 5,848 query rows, one primary hard mixed public stream, and reports 0 validation errors.

\subsection{Workload Axes}

\bench\ varies several workload axes simultaneously so that results can be interpreted diagnostically rather than only as a single aggregate score. The main axes are query family, predicate family, episode horizon, stream regime, temporal position, and evidence status.
The \emph{query-family} axis determines the answer shape: point values, bitemporal state, provenance evidence, change sets, uncertainty sets, preconditions, failure localization, or goal-state answers. The \emph{predicate-family} axis determines whether the maintained state involves unary object states, binary relations, numeric values, material states, or contact configurations. The \emph{episode-horizon} axis controls stream length and the number of opportunities for stale evidence, late repair, and accumulated uncertainty. The \emph{stream-regime} axis distinguishes clean streams from perturbed streams. The \emph{temporal-position} axis distinguishes initial, transition-near, and final queries. The \emph{evidence-status} axis captures whether the public stream provides support evidence, contradict evidence, both, or neither.
These axes determine how results should be read. A system that performs well on clean final-state checks may still fail on transition-near as-of queries. A system that handles Boolean unary predicates may not handle numeric temperature states or relation changes. A system that answers short-episode goal queries correctly may degrade on longer episodes because stale evidence accumulates. \bench\ therefore reports overall metrics together with query-family and stream-regime breakdowns.

\subsection{Perturbation Regimes}

The perturbation regimes are designed to stress data-management behavior rather than to model one specific sensor. Clean streams test basic state-view construction. Missing-observation regimes test unknown-state and stale-state handling. Temporal-disorder regimes test valid-time and transaction-time repair. Low-confidence regimes test confidence propagation and calibration. Conflict regimes test support/contradict evidence handling and status semantics.
The reported hard mixed stream combines several perturbation types and serves as the primary stress condition. This choice reflects the benchmark target: embodied observation streams rarely contain a single isolated defect. A system must often handle missing evidence, delayed evidence, low-confidence evidence, and contradictory evidence in the same episode. The perturbations are target-aware in the sense that they stress evidence near task-relevant states, while preserving enough background observations for realistic stream structure.

Perturbation provenance is kept hidden from systems under test. Public streams contain the resulting observations, but they do not reveal which records were removed, downgraded, delayed, or injected. Hidden perturbation metadata is retained only for diagnostic analysis and validation. This prevents the benchmark from becoming a label-reading task while still allowing error analysis by perturbation type.

\subsection{Evaluation Modes}

\bench\ supports both online and batch evaluation modes. In online mode, a system ingests observations in arrival-time order and answers queries at their declared transaction times. This mode best matches task monitors that operate during execution and emphasizes update cost, state repair, and incremental view maintenance.
In batch mode, a system may preprocess the full public stream before answering queries. However, it must still respect transaction-time constraints for bitemporal queries such as \texttt{AS\_OF\_STATE}. Batch mode is useful for comparing database implementations, retrieval systems, and offline analysis pipelines whose natural interface is not an online monitor.
Both modes are valid, but they answer different performance questions. Online mode evaluates whether a system can maintain and repair task-state views incrementally. Batch mode evaluates whether a system can implement the benchmark semantics correctly over the public stream. When both modes are available for a system, \bench\ reports preprocessing time separately from update time and query time, so that an offline full-stream scan is not confused with an online maintained view.

\subsection{Evaluation Metrics}

\bench\ reports diagnostic metrics rather than a single leaderboard number. Exact accuracy checks whether the predicted structured answer matches the hidden answer for a query. Value accuracy evaluates the value field when applicable. Status accuracy checks whether the system distinguishes \texttt{known}, \texttt{unknown}, \texttt{uncertain}, \texttt{conflict}, and \texttt{n/a}. Coverage reports how often the system returns a valid answer in the required schema.
Structured workloads use additional metrics. \texttt{STATE\_DIFF} uses both exact change-set match and mean F1 over changed state atoms. \texttt{CHECK\_GOAL} reports task-level correctness and goal-predicate F1, so that a system can receive partial credit when it identifies some but not all satisfied or violated goal predicates. Provenance-oriented workloads can score overlap between emitted support or contradict evidence and oracle evidence. For systems that emit confidence, \bench\ also reports confidence mean absolute error where hidden confidence is defined.
These metrics are intentionally diagnostic. A single global score would hide differences between stale-state errors, transaction-time errors, status-calibration errors, missing provenance, and change-set errors. \bench\ therefore reports overall metrics together with query-family, stream-regime, and structured-answer breakdowns. This design lets the benchmark identify not only which system performs better, but also which data-management capability is responsible for the difference.

